\section{Worked Example}
\label{app:worked-example}

% Wu & Tu's Appendix F holds case studies of the dependency structures their model infers. Ours
% holds the numeric example, which serves the same purpose: showing the reader what the model
% actually computes rather than asserting that it computes it.

\TODO{reproduce \texttt{causal\_pt\_output\_note.pdf} \S5 --- $\Vocab = \{\textit{the},
\textit{cat}, \textit{sat}, \textit{mat}\}$, $\nlab = 2$, one channel, both readout modes, with
printed posteriors and logits at every step.}

\NOTE{take the numbers from the test-suite output rather than retyping them, so the appendix and
the tests cannot drift apart, and say in the caption that they agree with the independently
derived reference.}

\TODO{if space allows, add the analogue of Wu \& Tu's Appendix F proper: a case study of the
dependency structures the causal model infers on real text. Their finding was that the induced
structures are only partially consistent with human intuition --- reporting the same honestly is
better than omitting the section.}
